\chapter{Lezione 9 --- Pre-warping, trasformata di Tustin e spettro discreto tramite FFT}

\section{Dal regolatore continuo \(R(s)\) al regolatore discreto \(\tilde R(z)\)}

L'obiettivo della lezione è capire come passare da un regolatore o da un filtro
progettato nel continuo a un corrispondente oggetto discreto.

Nel continuo si lavora con
\[
R(s),
\]
mentre nel discreto si vuole ottenere una funzione di trasferimento
\[
\tilde R(z).
\]

Il legame esatto tra le due descrizioni è dato dalla relazione
\[
z=e^{sT},
\qquad\text{equivalentemente}\qquad
s=\frac{1}{T}\ln z.
\]

Tuttavia questa relazione è trascendente e non è comoda da usare direttamente
per ottenere una funzione razionale in \(z\).
Per questo si usa una sostituzione approssimata, ma razionale, cioè la
\textbf{trasformata di Tustin}.

\section{Trasformata di Tustin}

La trasformata di Tustin, o trasformata bilineare, si scrive come
\[
s=\frac{2}{T}\frac{z-1}{z+1}.
\]

In forma inversa, si può anche scrivere
\[
z=\frac{\frac{2}{T}+s}{\frac{2}{T}-s}.
\]

Questa trasformazione ha due proprietà molto importanti:
\begin{itemize}
    \item manda il semipiano sinistro del piano \(s\) all'interno del cerchio unitario;
    \item manda l'asse immaginario del piano \(s\) sul cerchio unitario del piano \(z\).
\end{itemize}

Quindi:
\begin{itemize}
    \item poli stabili nel continuo vengono mappati in poli stabili nel discreto;
    \item la stabilità viene preservata.
\end{itemize}

\section{Risposta armonica e problema della distorsione}

Se si discretizza un filtro analogico con Tustin, la risposta armonica del filtro discreto
non coincide esattamente con quella del filtro continuo: la scala delle frequenze viene distorta.

Per un filtro passa-basso, in genere interessa preservare bene il comportamento
in una frequenza significativa, tipicamente la frequenza di taglio.

Ad esempio, se nel continuo si ha un filtro del primo ordine,
alla pulsazione di taglio \(\omega_b=1/\tau\) si ha tipicamente:
\begin{itemize}
    \item modulo pari a \(-3\) dB rispetto al valore a bassa frequenza;
    \item fase pari a \(-45^\circ\).
\end{itemize}

L'idea è allora imporre che queste caratteristiche siano preservate,
non ovunque, ma almeno in una pulsazione di interesse.

\section{Pre-warping}

Per correggere la distorsione della scala in frequenza si introduce una costante \(k\)
nella trasformata bilineare:
\[
s=k\frac{z-1}{z+1}.
\]

Si sceglie \(k\) in modo che la corrispondenza tra filtro analogico e filtro discreto
sia esatta in una pulsazione \(\omega_0\) prefissata.

Il valore corretto di \(k\) è
\[
k=\frac{\omega_0}{\tan\!\left(\frac{\omega_0 T}{2}\right)}.
\]

Questa scelta prende il nome di \textbf{pre-warping} alla pulsazione \(\omega_0\).

\subsection*{Condizione di progetto}

Con questa scelta si forza l'uguaglianza
\[
R(j\omega_0)=\tilde R(e^{j\omega_0 T})
\]
nel senso della risposta armonica: il filtro discreto riproduce esattamente
il comportamento del filtro continuo alla frequenza \(\omega_0\).

\subsection*{Caso limite a bassa frequenza}

Se \(\omega_0\to 0\), si ha
\[
\tan\!\left(\frac{\omega_0 T}{2}\right)\sim \frac{\omega_0 T}{2},
\]
e quindi
\[
k=\frac{\omega_0}{\tan(\omega_0 T/2)}
\longrightarrow
\frac{\omega_0}{\omega_0 T/2}
=
\frac{2}{T}.
\]

Quindi la trasformata di Tustin standard è il caso limite del pre-warping
quando si vuole privilegiare la bassa frequenza.

\section{Più frequenze di interesse}

Negli appunti compare un'osservazione importante:
se ci sono più frequenze da preservare, non è possibile renderne esatta la corrispondenza
per tutte contemporaneamente con una sola trasformazione di Tustin pre-warpata.

In pratica:
\begin{itemize}
    \item si sceglie una frequenza prioritaria;
    \item oppure si lavora per tentativi e compromessi;
    \item spesso si privilegia la frequenza di taglio di un passa-basso.
\end{itemize}

Questo è coerente con il fatto che la trasformazione resta un'approssimazione,
non una equivalenza esatta tra continuo e discreto.

\section{Campionamento di una sequenza lunga e calcolo dello spettro}

Nella seconda parte della lezione si passa al problema pratico:
se si campiona una lunga sequenza di dati e si vuole ricavare lo spettro,
si calcola la trasformata di Fourier discreta della sequenza campionata.

Se il segnale continuo \(x(t)\) viene campionato con periodo \(T\), si ottiene la successione
\[
\tilde x(kT),\qquad k=0,1,\dots,N-1.
\]

A partire da questi \(N\) campioni si costruisce il vettore dei campioni nel tempo
e si calcola la trasformata discreta di Fourier:
\[
\tilde X(h)=\sum_{k=0}^{N-1}\tilde x(kT)\,e^{-j\frac{2\pi}{N}hk},
\qquad
h=0,1,\dots,N-1.
\]

Il risultato \(\tilde X(h)\) è un vettore di numeri complessi che rappresenta modulo e fase
dello spettro campionato.

\section{FFT}

Il calcolo diretto della DFT richiederebbe un numero di operazioni dell'ordine di
\[
N^2.
\]

Per rendere il calcolo efficiente si usa la \textbf{FFT} (\emph{Fast Fourier Transform}),
cioè un algoritmo che calcola la DFT con costo molto minore, tipicamente dell'ordine di
\[
N\log N.
\]

Per questo motivo, quando si devono elaborare sequenze lunghe,
si preferisce usare la FFT invece della formula diretta della DFT.

\subsection*{Osservazione pratica}

È molto conveniente scegliere un numero di campioni \(N\) che sia una potenza di \(2\),
perché molti algoritmi FFT lavorano in modo particolarmente efficiente in questo caso.

\section{Asse delle frequenze restituito dalla FFT}

Un punto essenziale della lezione è capire che i valori \(\tilde X(h)\)
non sono associati a frequenze arbitrarie, ma a frequenze equispaziate.

La pulsazione corrispondente al bin \(h\) è
\[
\omega_h = h\,\frac{2\pi}{NT}.
\]

Quindi:
\[
\Delta\omega = \frac{2\pi}{NT}
\]
è il passo tra due campioni consecutivi in frequenza.

In particolare:
\[
h=0 \quad \Longrightarrow \quad \omega_0=0,
\]
\[
h=\frac{N}{2} \quad \Longrightarrow \quad \omega_{N/2}=\frac{\pi}{T}.
\]

Quindi il bin \(h=N/2\) corrisponde alla pulsazione di Nyquist.

\subsection*{Osservazione sull'intervallo utile}

Per segnali reali, lo spettro contiene informazione ridondante nella seconda metà,
per simmetria complessa coniugata.
Per questo, in molti casi pratici, l'intervallo utile è solo
\[
0\le \omega \le \frac{\pi}{T},
\]
cioè da \(0\) fino alla metà della pulsazione di campionamento.

In questo intervallo si hanno, a seconda della convenzione di indicizzazione,
circa \(N/2\) oppure \(N/2+1\) campioni utili.

\section{Interpretazione MATLAB}

Negli appunti si osserva che ambienti come MATLAB, quando si usa \texttt{fft},
restituiscono il vettore dei bin senza costruire automaticamente la scala fisica delle frequenze.

Di fatto, se non si imposta esplicitamente il periodo di campionamento \(T\),
la rappresentazione sembra fatta come se \(T=1\).

Quindi, per riportare correttamente l'asse delle pulsazioni,
bisogna costruire a mano il vettore
\[
\omega_h = h\,\frac{2\pi}{NT}.
\]

\section{Scala lineare e scala logaritmica}

Negli appunti viene poi discusso il problema della rappresentazione grafica dello spettro.

I campioni in frequenza prodotti dalla FFT sono \textbf{equispaziati in scala lineare},
cioè il passo tra una frequenza e la successiva è costante:
\[
\Delta\omega = \frac{2\pi}{NT}.
\]

Se però lo spettro viene riportato su asse logaritmico, i punti non appaiono più equidistanti.

In particolare:
\begin{itemize}
    \item alle basse frequenze i punti si distinguono bene;
    \item alle alte frequenze i punti si affollano e si distinguono sempre meno.
\end{itemize}

Questo è un fatto puramente geometrico dovuto al cambio di scala,
non a una diversa densità dei campioni della FFT.

\section{Limite inferiore della griglia spettrale}

La minima pulsazione positiva rappresentabile con \(N\) campioni e periodo \(T\) è
\[
\omega_{\min}=\frac{2\pi}{NT}.
\]

Questa è la frequenza del primo bin non nullo.
Per questo, se si vuole vedere bene il comportamento a bassissima frequenza,
bisogna aumentare \(N\), cioè osservare il segnale per un tempo più lungo.

\section{Normalizzazione dell'ampiezza}

Negli appunti compare un'ultima osservazione importante:
nel passaggio dal dominio del tempo alla trasformata discreta di Fourier
possono comparire discrepanze nei valori delle ampiezze se non si usa la corretta normalizzazione.

I coefficienti di scala dipendono dalla convenzione scelta:
\begin{itemize}
    \item definizione di trasformata continua;
    \item definizione di DFT/FFT;
    \item spettro monolatero o bilatero;
    \item ampiezza di picco oppure valore efficace (RMS).
\end{itemize}

Per questo, in applicazioni pratiche, possono comparire fattori del tipo:
\[
\frac{1}{N},\qquad
T,\qquad
2,\qquad
\frac{1}{2\pi},\qquad
\sqrt{2},
\]
o combinazioni di essi, a seconda di ciò che si vuole rappresentare.

\subsection*{Interpretazione}

Se si vuole ottenere dallo spettro FFT un'ampiezza fisicamente coerente con il segnale originale,
bisogna sempre specificare:
\begin{itemize}
    \item quale trasformata si sta usando;
    \item come si sta costruendo l'asse delle frequenze;
    \item quale normalizzazione si sta adottando.
\end{itemize}

\section{Modulo in dB}

Quando si vuole rappresentare il modulo dello spettro in forma logaritmica,
si usa tipicamente
\[
20\log_{10}|\tilde X(h)|.
\]

In questo modo si ottiene il diagramma del modulo in dB,
analogo a quello usato nei diagrammi di Bode.

\section{Sintesi della lezione}

I punti principali della lezione sono:
\begin{enumerate}
    \item il passaggio \(R(s)\mapsto \tilde R(z)\) si fa in pratica tramite la trasformata di Tustin;
    \item la trasformata di Tustin standard è
    \[
    s=\frac{2}{T}\frac{z-1}{z+1};
    \]
    \item con il pre-warping si usa
    \[
    s=k\frac{z-1}{z+1},
    \qquad
    k=\frac{\omega_0}{\tan(\omega_0 T/2)},
    \]
    per preservare la risposta armonica alla frequenza \(\omega_0\);
    \item al limite \(\omega_0\to 0\) si ritrova
    \[
    k\to \frac{2}{T};
    \]
    \item una sequenza di \(N\) campioni può essere trasformata nel dominio della frequenza
    tramite DFT o, più efficientemente, tramite FFT;
    \item il costo computazionale della FFT è molto più basso di quello della DFT diretta;
    \item i bin spettrali sono equispaziati in scala lineare con passo
    \[
    \Delta\omega=\frac{2\pi}{NT};
    \]
    \item la pulsazione del bin \(h\) è
    \[
    \omega_h=h\frac{2\pi}{NT};
    \]
    \item per segnali reali, l'intervallo utile è
    \[
    0\le \omega \le \frac{\pi}{T};
    \]
    \item su asse logaritmico i punti FFT risultano sempre più fitti alle alte frequenze;
    \item per rappresentare correttamente le ampiezze occorre usare la normalizzazione coerente
    con la convenzione adottata.
\end{enumerate}